% !TEX root = ../Abschlussarbeit_DE.tex
\section{Modulare Softwarearchitektur mittels Module Federation}
Jede Großbaustelle hat ihre eigenen individuellen Herausforderungen. Da jedoch viele Softwarelösungen als Allround-Lösungen konzipiert sind, 
scheitern sie oft an den spezifischen Detailanforderungen oder sind zu komplex und können so ihr Potenzial bei Weitem nicht ausschöpfen \cite[Kap.~3.1]{feth2023studie}. 
In der Praxis führt dies oft zu individuellen Spezialaufträgen an die Softwarehersteller - ein Prozess, der meist mit erheblichen Zusatzkosten und langen Entwicklungszeiten 
verbunden ist \cite[Kap.~4.1]{feth2023studie}. 
Um dem entgegenzuwirken, bedarf es einer flexiblen und vor allem einfach erweiterbaren Struktur, in welcher individuelle Softwarelösungen unkompliziert implementiert werden können. 
Ein möglicher Ansatzpunkt hierfür ist das Konzept der Microfrontends. Dabei werden einzelne Funktionalitäten als eigenständige Webapplikationen entwickelt, 
die sich zu einer gemeinsamen Benutzeroberfläche zusammenfügen lassen. 
%% Einzelne Gewerke können einfach hinzugefügt werden oder erp software
%% Vorteile gegenüber Monolithen (z.B. unabhängige Entwicklung und Deployment, skaliebarkeit, spezielle Lösungen)
%% Wichtige Dinge wie Kommunikation erwähnen 
\subsection{Aufbau von Microfrontends}
%% Aufbau (einzelne Komponenten wie z.B. Header )
%% Erklärung von Shell und Routing
\subsection{Kommunikation zwischen Microfrontends}
%% LocalStorage oder andere muss noch recherchiert ewrden
%% Nutzen von GraphQL mit listeners
\subsection{Konfiguration einer Microfrontendarchitektur}
%% Beispiel Vite
%% Erwähnung Monorepo
%% relativ straight forward wie gemacht wird
%% Shell Scripte
%% Konfigdatei für ports etc ver microfrontends